% SPDX-License-Identifier: CC-BY-4.0
% Copyright 2025 Florian Aram Feuerriegel — kassensturz.org
% Abschnitte 6 und 7 — Grenzen des Modells, Schluss

\section{Grenzen}
\label{sec:grenzen}

\paragraph{Struktur.} Das Modell kennt kein allgemeines Gleichgewicht: Löhne,
Preise und Zinsen reagieren nicht auf die Politik, der Zins ändert sich nur
im Zinsschock. Innerhalb eines Haushaltstyps sind alle gleich, außer im
obersten Prozent und bei der Armutsquote; ein Typ nahe einer Tarifgrenze
reagiert deshalb sprunghaft, wie D3 beim Grundfreibetrag. Verhalten wirkt über
je eine begrenzte Elastizität, ohne Erwartungen und ohne Entscheidung über die
Erwerbsbeteiligung; Kapitaleinkünfte reagieren gar nicht, weil für sie keine
belegte Elastizität vorliegt. Das mittlere Arbeitsangebot $\bar\lambda$, das
das BIP der Folgeperiode bewegt, ist mit der Zahl der Haushalte gewichtet, nicht
mit ihren Löhnen. Eine Reaktion in D3 zählt so viel wie eine in D9, und beim
Grundfreibetrag, dessen Angebotswirkung vor allem aus D3 kommt, fällt die
Wirkung auf das BIP-Niveau fast doppelt so hoch aus wie mit Lohngewichten. Die
Nachfrage hebt BIP und Einnahmen, aber nicht die
Einkommen der Haushaltstypen. Die Buchungsidentität gilt damit für die
direkte Wirkung einer Politik, nicht für ihre zweite Runde, und der
Multiplikator hängt weder von der Konjunktur noch von der Geldpolitik ab.

\paragraph{Buchung.} Der Basisfaktor der Umsatzsteuer, $\varphi = 1{,}68$,
skaliert beim Staat auch Änderungen des Haushaltskonsums, bei den Haushalten
nur Satzänderungen. Eine Einkommensänderung bringt dem Staat daher rund 60\,\%
mehr Umsatzsteuer, als die Haushalte zahlen. Auch bei Satzänderungen rechnen
beide Seiten mit verschiedenen Bemessungsgrundlagen, der Staat mit Lohnsumme und
Konsumreaktion, die Haushalte mit den Arbeitseinkommen und dem Konsum ihrer
Typen. Bei den gemessenen Schritten von Umsatzsteuer, Beitragssätzen und
Beitragsbemessungsgrenze weichen sie um bis zu ein Zehntel voneinander ab, und
die Nachfragewirkung folgt der Seite der Haushalte. Die Beitragsbemessungsgrenze
wirkt über eine Faustregel ohne Primärtabelle: Eine um $x$\,\% andere Grenze
ändert die Lohnsumme um $0{,}12\,x$\,\%, gleich wie weit sie sich bewegt, und
weder private Krankenversicherung noch Beamte sind abgetrennt. Die Vermögensteuer greift erst ab 2~Mio.\,\euro, ihre
Inzidenz folgt aber dem gesamten Vermögen: Ein Teil der Last landet bei
Haushaltstypen, die sie nicht zahlen würden. Die Rentenzahlungen der
Haushalte (362~Mrd.~\euro) und der Demografieaufschlag (auf 390~Mrd.~\euro)
beruhen auf verschiedenen Abgrenzungen der Rentenausgaben.

\paragraph{Daten.} Mehrere Größen sind Näherungen ohne Primärtabelle, weil die
Quellen beim Anlegen nicht erreichbar waren: die Rentenanteile je Dezil, das
Profil der Grenzkonsumneigung, der öffentliche Kapitalstock, das
Milliardärsvermögen und der Pareto-Exponent; für die Wegzugsreaktion im
obersten Prozent fehlt ein Beleg. Der Spitzensatz trifft über den Pareto-Rand
nur wenig Einkommen, ein Prozentpunkt bringt mechanisch
\Wb{zerl}{spitze}{mechanisch}~Mrd.~\euro; ob der Rand mit $\alpha = 1{,}5$ die
Einkommen oberhalb der Spitzenzone trifft, ist an der Einkommensteuerstatistik
zu prüfen. Die Ausgabenseite ist grob gegliedert, und
der Ausgleichsposten nimmt Fehler auf, die ein Abgleich Posten für Posten mit
der VGR finden würde. Die Beitragseinnahmen liegen um ein Achtel zu hoch
(Abschnitt~\ref{sec:kalibrierung}). Nur die Rente altert; Gesundheit, Pflege
und das schrumpfende Erwerbspersonenpotenzial fehlen, und die Demografie folgt
der 14.~statt der aktuellen 15.~koordinierten Bevölkerungsvorausberechnung.

\paragraph{Rahmen.} Die Schuldenbremse gilt für den Gesamtstaat, nicht
getrennt für Bund und Länder; die europäischen Fiskalregeln fehlen. Zwei der
fünf Schockwirkungen, Investitionsrückgang und Emissionsminderung, stehen in
der Bibliothek, werden aber nicht gerechnet. Einen Klimaschaden auf das
deutsche BIP gibt es bewusst nicht: Deutschland verursacht rund 1,5\,\% der
weltweiten Emissionen, der Rückkanal auf die eigene Wirtschaft ist
vernachlässigbar. Das Klimaziel steht über Emissionen und Budget im Modell,
nicht über das BIP.

\section{Schluss}
\label{sec:schluss}

Ein Modell mit zwölf Haushaltstypen, Elastizitäten aus der Literatur und einem
einfachen Makrorahmen bildet die zentralen Zielkonflikte der deutschen
Finanzpolitik ab, wenn zwei Bedingungen erfüllt sind: Staat und Haushalte
buchen dieselben Beträge, und Tests prüfen Eigenschaften statt Ist-Werte. Unter
diesen Bedingungen hat keine Stellgröße eine Wirkung ohne Preis. Einnahmen
kosten Gleichheit oder Wachstum, Entlastungen kosten Saldo, und wer wie viel
trägt, entscheidet über Größe und Richtung der Nebenwirkungen.

Die Weiterentwicklung folgt den Grenzen: Primärtabellen statt Näherungen für
Renten, Konsumneigung und Beitragsbemessungsgrenze, ein Arbeitsangebot mit
Lohngewichten, eine Nachfrage, die auf die Einkommen der Haushalte zurückwirkt,
die Inzidenz der Vermögensteuer nach dem Freibetrag, die Trennung von Bund und
Ländern mit den europäischen Fiskalregeln und eine Demografie, die über die
Rente hinausreicht.
